\documentclass[11pt]{article}

\usepackage{amsmath,amssymb,amsfonts,amsthm}
\usepackage[margin=1in]{geometry}
\usepackage{hyperref}
\usepackage{booktabs}
\usepackage{enumitem}
\usepackage{graphicx}
\usepackage{xcolor}
\usepackage{mathrsfs}
\usepackage[expansion=false]{microtype}

% Theorem environments
\newtheorem{theorem}{Theorem}[section]
\newtheorem{lemma}[theorem]{Lemma}
\newtheorem{proposition}[theorem]{Proposition}
\newtheorem{corollary}[theorem]{Corollary}
\theoremstyle{definition}
\newtheorem{definition}[theorem]{Definition}
\newtheorem{assumption}[theorem]{Assumption}
\theoremstyle{remark}
\newtheorem{remark}[theorem]{Remark}

% Custom commands consistent with the geometric-economics book
\newcommand{\E}{\mathcal{E}}
\newcommand{\BF}{\mathrm{BF}}
\newcommand{\R}{\mathbb{R}}
\newcommand{\BGE}{\mathrm{BGE}}
\newcommand{\Nash}{\mathrm{NE}}
\DeclareMathOperator*{\argmin}{arg\,min}
\DeclareMathOperator*{\argmax}{arg\,max}

\title{\textbf{The Computational Origin of Bounded Rationality:\\
A Tractability Theorem for Behavioral Equilibria on the Economic Decision Manifold}}

\author{Andrew~H.~Bond\\
Department of Computer Engineering\\
San Jos\'{e} State University\\
\texttt{andrew.bond@sjsu.edu}\\
ORCID: 0009-0003-2599-6158}

\date{June 2026}

\begin{document}
\maketitle

\begin{abstract}
\noindent
Computing a Nash equilibrium in a general bimatrix game is PPAD-complete, and the intractability worsens
for continuous-action and multi-agent settings. The Bond Geodesic Equilibrium (BGE), a generalization
of Nash to multi-dimensional decision manifolds, is therefore also PPAD-hard in general. Yet real
agents---human and artificial---routinely play behavioral equilibria that classical theory says they
cannot compute. We resolve the apparent contradiction by giving the formal complexity-theoretic
content of bounded rationality.

We prove that a Bond Geodesic Equilibrium becomes computationally tractable when the underlying
decision manifold admits an \emph{admissible heuristic truncation}---a restriction to a small active
subset of dimensions on which an $A^*$-style search with an admissible heuristic locates a
near-optimal path. Specifically, under (i) admissibility on the truncated manifold and (ii) a
contraction condition on the cross-agent coupling, the unique truncated BGE is computable in time
polynomial in the number of agents and exponential only in the truncation dimension, parameterized
by the heuristic--coupling ratio. The PPAD-hardness is preserved on the full single-action manifold-game
class; it is bypassed on the contractive truncated subclass that an actual agent can search, in
the sense that contraction plus admissible truncation jointly place the equilibrium problem in
polynomial time.

The reframing is constructive: bounded rationality is not a psychological deficit but the
computational price tractable behavior must pay on an otherwise intractable equilibrium landscape.
The framework's interpretable dimensions are not merely descriptive labels for choice frequencies
but candidate heuristic coordinates over which admissible truncation is performed.

We anchor the theorem empirically using a recently published cross-domain validation that calibrates
the 9-dimensional decision manifold on bargaining and public-goods aggregate data and tests the
resulting metric on prospect-theory phenomena and high-stakes ultimatum data. The empirically
selected active set has dimension $k=3$, and the tractability theorem predicts the observed
equilibrium structure to within the validated tolerance.
\end{abstract}

\medskip
\noindent\textbf{Keywords:} bounded rationality, PPAD-completeness, Nash equilibrium, Bond Geodesic
Equilibrium, admissible heuristic, $A^*$ search, behavioral game theory, computational complexity,
geometric economics.

\medskip
\noindent\textbf{JEL classification:} C72 (Noncooperative Games), C73 (Stochastic and Dynamic Games),
D03 (Behavioral Economics), D81 (Criteria for Decision-Making under Risk and Uncertainty).

\noindent\textbf{ACM classification:} F.2.0 (Analysis of Algorithms and Problem Complexity, General);
J.4 (Computer Applications: Social and Behavioral Sciences).


%==============================================================================
\section{Introduction}
\label{sec:intro}
%==============================================================================

\subsection{The complexity gap in equilibrium theory}

Classical game theory rests on three pillars: existence of Nash equilibrium~\cite{nash1950}, its
interpretation as a self-consistent prediction of strategic behavior, and the implicit assumption
that real agents can compute or learn it. The first pillar is mathematically secure. The third has
been undermined by complexity theory.

\cite{daskalakis2009} proved that computing a Nash equilibrium is \textsc{PPAD}-complete, and
\cite{chen2009} sharpened this to two-player games. The conclusion is not that equilibria fail to
exist---they do, by Nash's theorem---but that no polynomial-time algorithm can find them under
standard complexity assumptions. The gap between existence and computability widens in continuous-action,
many-player, and incomplete-information settings, where even approximation becomes hard. Real agents,
human or artificial, cannot be performing the global optimization that equilibrium analysis
implicitly attributes to them.

Behavioral economics has documented the consequence for half a century. From Simon's
satisficing~\cite{simon1955} through Kahneman and Tversky's heuristics-and-biases
programme~\cite{kahneman1979,tversky1992} and the modern resource-rationality
framework~\cite{lieder2020}, observed behavior departs systematically from the equilibrium
predictions of scalar utility maximization. The standard interpretation is that humans are
\emph{boundedly rational}: they use heuristics because they are cognitively limited.

This paper develops the converse interpretation. Heuristics are not concessions to cognitive limits;
they are the computational mechanism by which an otherwise intractable equilibrium becomes tractable
on the restricted class behavioral data actually populate. The PPAD barrier is not defeated in any
broad sense, but it is bypassed where it matters most: the standard equilibrium question demands
fixed points on the full action space, when an admissible truncation paired with a contraction
condition suffices, and the empirically observed behavioral class satisfies both.

\subsection{The Bond Geodesic Equilibrium}

The geometric framework~\cite{bond2026geometric,bond2026tcss} models economic decisions as
movement on a 9-dimensional decision manifold $\E$ whose dimensions encode monetary consequences,
rights, fairness, autonomy, privacy/trust, social impact, virtue/identity, legitimacy, and epistemic
status. A diagonal covariance metric $\Sigma$ assigns Mahalanobis cost to displacement; an agent's
optimal action is the minimum-cost path---the \emph{Bond geodesic}---from their starting state to a
goal region. When multiple agents interact, the equilibrium concept is the \emph{Bond Geodesic
Equilibrium} (BGE): each agent's chosen path is optimal given the paths of the others. The BGE
generalizes Nash equilibrium: when all non-monetary dimensions are inactive, BGE reduces to Nash on
the scalar projection.

The framework has been validated empirically by~\cite{bond2026tcss}, which calibrates the metric on
nine game-theoretic targets (ultimatum and dictator games, public-goods rounds), tests it without
re-tuning on six prospect-theory targets and one historical replication, achieves 16/16 matches at
MAE 2.70\%, and falsifies the framework's stake-scaling-invariance corollary on
Andersen et al.'s~\cite{andersen2011} high-stakes Indonesia ultimatum data ($n{=}458$) in exactly the
direction the framework predicts. The empirically selected active set has dimension $k=3$:
$\{d_6, d_7, d_9\}$ (Social Impact, Virtue/Identity, Epistemic Status).

\subsection{Contribution}

This paper provides the formal complexity-theoretic content of bounded rationality in the geometric
framework. The contributions are:

\begin{enumerate}
\item \textbf{A PPAD-hardness result for BGE} (Theorem~\ref{thm:bge-ppad}). Computing an exact BGE is
\textsc{PPAD}-complete in general, by reduction from Nash equilibrium via the BGE--Nash equivalence
established in~\cite{bond2026geometric}.

\item \textbf{A tractability theorem under admissible heuristic truncation}
(Theorem~\ref{thm:tractable-bge}). When the manifold admits an admissible heuristic truncation to a
$k$-dimensional active set and the cross-agent coupling satisfies a contraction condition, the unique
truncated BGE is computable in time polynomial in the number of agents and exponential only in $k$,
parameterized by the heuristic--coupling ratio.

\item \textbf{A reframing of bounded rationality} (Section~\ref{sec:reframe}). Heuristic truncation
is not a deviation from rational behavior but the computational price tractable equilibrium must pay.
Real agents play approximate BGE because the truncated landscape admits a polynomial-time best-response
contraction; they cannot play exact BGE because the full landscape is \textsc{PPAD}-hard.

\item \textbf{An empirical anchor} (Section~\ref{sec:empirical}). The validated TCSS calibration is
precisely the instantiation of admissible heuristic truncation with $k=3$ that the tractability
theorem requires. We map the published empirical parameters to the theorem's preconditions and verify
that the parameters satisfy them.
\end{enumerate}

\subsection{Organization}

Section~\ref{sec:prelim} fixes notation, briefly recalls the geometric framework, and states the
results of~\cite{bond2026geometric,bond2026tcss} this paper builds on without re-deriving.
Section~\ref{sec:complexity} establishes the PPAD-hardness of general BGE. Section~\ref{sec:truncation}
defines admissible heuristic truncation and proves the main tractability theorem.
Section~\ref{sec:empirical} provides the empirical anchor. Section~\ref{sec:reframe} discusses what
the result implies for the theoretical status of bounded rationality. Section~\ref{sec:limitations}
identifies limitations and open questions. Section~\ref{sec:conclusion} concludes.


%==============================================================================
\section{Preliminaries}
\label{sec:prelim}
%==============================================================================

We collect the definitions and prior results required for the complexity arguments. Throughout, we
adopt the notation of~\cite{bond2026geometric}; readers familiar with that monograph may skim this
section. None of the theorems in this section are new.

\subsection{The economic decision complex}

\begin{definition}[Economic decision complex, \cite{bond2026geometric}]
\label{def:complex}
The \emph{economic decision complex} $\E = (V, E, w)$ is a finite weighted directed graph whose
vertices $V \subset \R^9$ are economic states embedded in the 9-dimensional decision space spanning
\textsf{consequences} ($d_1$), \textsf{rights} ($d_2$), \textsf{fairness} ($d_3$), \textsf{autonomy}
($d_4$), \textsf{privacy/trust} ($d_5$), \textsf{social impact} ($d_6$), \textsf{virtue/identity}
($d_7$), \textsf{legitimacy} ($d_8$), and \textsf{epistemic status} ($d_9$). Each directed edge
$(v_i, v_j) \in E$ represents an admissible action with weight
\begin{equation}
\label{eq:edge-weight}
w(v_i, v_j) \;=\; \sqrt{\Delta\mathbf{a}^\top \Sigma^{-1} \Delta\mathbf{a}}
                \;+\; \sum_{k \in \mathcal{B}(v_i, v_j)} \beta_k,
\end{equation}
where $\Delta\mathbf{a} = v_j - v_i$, $\Sigma$ is a $9 \times 9$ positive-definite covariance matrix,
and $\mathcal{B}(v_i, v_j)$ is the set of moral boundaries (with penalties $\beta_k \geq 0$) crossed
by the transition $v_i \to v_j$.
\end{definition}

We assume $\Sigma$ is diagonal throughout, $\Sigma = \mathrm{diag}(\sigma_1^2, \ldots, \sigma_9^2)$;
the non-diagonal case is treated in~\cite{bond2026geometric} and not required here. We allow
$\sigma_k^2 = \infty$, in which case dimension $k$ is \emph{inactive} (movement on $d_k$ contributes
zero cost).

\subsection{Bond geodesics and the BGE}

A \emph{Bond geodesic} from $v_0 \in V$ to a goal region $G \subseteq V$ is a minimum-cost path
$\gamma^* = \argmin_\gamma \sum_{i} w(v_i, v_{i+1})$ subject to $\gamma(0) = v_0$ and
$\gamma(\mathrm{end}) \in G$. A geodesic is computed in time $O(|E| + |V|\log|V|)$ by Dijkstra
or $A^*$ search on $\E$.

\begin{definition}[Bond Geodesic Equilibrium, \cite{bond2026geometric}]
\label{def:bge}
A profile $(\gamma_1^*, \ldots, \gamma_n^*)$ of paths on agent-indexed decision complexes
$\{\E_i\}_{i=1}^n$ is a \emph{Bond Geodesic Equilibrium} if for every agent $i$, the path
$\gamma_i^*$ minimizes agent $i$'s total path cost
$\BF_i(\gamma_i \mid \gamma_{-i})$ over the set $S_i$ of simple paths from $v_0^{(i)}$ to $G_i$ in
$\E_i$, given the strategies $\gamma_{-i}^* = (\gamma_j^*)_{j \neq i}$ of the other agents.
\end{definition}

The \emph{augmented game}~\cite[Definition~20]{bond2026geometric} associated with a manifold game
$\Gamma = (N, \{\E_i\}, \{G_i\})$ is the normal-form game $\Gamma^+ = (N, \{S_i\}, \{u_i\})$ with
$u_i(\gamma_i, \gamma_{-i}) = -\BF_i(\gamma_i \mid \gamma_{-i})$. A direct calculation
\cite[Theorem~21]{bond2026geometric} yields the equivalence we will use repeatedly:

\begin{theorem}[BGE--Nash equivalence, \cite{bond2026geometric}]
\label{thm:equivalence}
A profile is a BGE of $\Gamma$ if and only if it is a Nash equilibrium of $\Gamma^+$. Existence of a
mixed BGE follows from Nash~(1950) applied to $\Gamma^+$.
\end{theorem}

This is the bridge between the manifold model and classical equilibrium theory. We will exploit it in
both directions: from $\Gamma^+$ to $\Gamma$ for hardness, from $\Gamma$ to $\Gamma^+$ for tractability.

\subsection{Best-response contraction}

\begin{lemma}[Contraction condition, \cite{bond2026geometric}]
\label{lem:contraction}
For each agent $i$, let $\alpha_i > 0$ be a lower bound on the self-separation of agent $i$'s
best-response mapping $T_i$ (how much $\BF_i$ changes with $i$'s own path), and let
$\kappa_i \geq 0$ be an upper bound on the cross-sensitivity (how much $\BF_i$ changes with the
others' paths). If $\alpha_i > \kappa_i$ for all $i \in N$, then the joint best-response
mapping $T : \prod_i \Delta(S_i) \to \prod_i \Delta(S_i)$ is a contraction with Lipschitz constant
$L = \max_i (\kappa_i/\alpha_i) < 1$. By Banach's theorem, the BGE is unique and iterated
best-response converges geometrically.
\end{lemma}

We refer to $L$ as the \emph{coupling ratio} of the game.


%==============================================================================
\section{The Computational Complexity of BGE}
\label{sec:complexity}
%==============================================================================

We first establish that without truncation, the BGE inherits the worst-case complexity of Nash. The
result is bounded in two ways below: we restrict the manifold game class to ensure polynomial
representation of the augmented game, and we supply an explicit gadget for the hardness reduction.
The combination is what a \textsc{PPAD}-completeness statement requires; the unrestricted manifold
game with exponentially many simple paths is a separate (and likely harder) problem we do not treat.

\begin{definition}[Single-action manifold game]
\label{def:single-action}
A manifold game $\Gamma$ is \emph{single-action} if for each agent $i$ the strategy set $S_i$
consists of two-edge paths $v_0^{(i)} \to v_{i,a} \to g^{(i)}$, one per pure action $a$ in a
polynomial-size action set $A_i$. The size of the input to a single-action manifold game is
$|N| + \sum_i |A_i| + |\Sigma|$, polynomially-bounded in $n$ and the per-agent action counts.
\end{definition}

\begin{theorem}[\textsc{PPAD}-completeness of single-action BGE]
\label{thm:bge-ppad}
Let $\varepsilon = 1/\mathrm{poly}(n)$. The problem of computing an $\varepsilon$-approximate BGE of a
single-action manifold game (Definition~\ref{def:single-action}) is \textsc{PPAD}-complete, even when
the covariance $\Sigma$ is diagonal with $\sigma_k^2 \in \{1, \infty\}$ for each $k$ and all boundary
penalties vanish ($\beta_k = 0$).
\end{theorem}

\begin{proof}
\textbf{Membership.} Definition~\ref{def:single-action} bounds $|S_i| = |A_i|$ polynomially in the
input size, so the augmented game $\Gamma^+$ has polynomial-size strategy sets and polynomial-size
payoff representation: $u_i(a, \mathbf{a}_{-i}) = -w(v_0^{(i)}, v_{i,a}) - w(v_{i,a}, g^{(i)})$ where
each $w$ is a single Mahalanobis evaluation computable in time $O(d)$ for dimension $d = 9$. By
Theorem~\ref{thm:equivalence}, every $\varepsilon$-BGE of $\Gamma$ is an $\varepsilon$-Nash
equilibrium of $\Gamma^+$, which is a finite normal-form game with polynomial-size representation. By
Daskalakis et al.~\cite{daskalakis2009}, computing $\varepsilon$-Nash in such games is in
\textsc{PPAD}.

\textbf{Hardness via explicit gadget.} Let $G = (\{1,2\}, \{A_1, A_2\}, \{p_1, p_2\})$ be a two-player
bimatrix game with payoff entries normalized to $[0,1]$ (rescaling does not change the
$\varepsilon$-Nash structure beyond a known affine factor; see~\cite{chen2009}). Computing
$\varepsilon$-Nash of $G$ for $\varepsilon = 1/\mathrm{poly}(\sum_i |A_i|)$ is \textsc{PPAD}-hard.

Construct a single-action manifold game $\Gamma$ as follows. For each agent $i \in \{1,2\}$ and each
action $a \in A_i$, let $v_{i,a} \in \R^9$ have coordinates
\[
(v_{i,a})_k \;=\; \begin{cases}
\sqrt{1 - p_i(a, \mathbf{a}_{-i})} & k = 1 \text{ (monetary)} \\
0 & k = 2, \ldots, 9
\end{cases}
\]
\emph{evaluated at the realized opponent action $\mathbf{a}_{-i}$.} The path
$v_0^{(i)} \to v_{i,a} \to g^{(i)}$ has cost
$\BF_i(a \mid \mathbf{a}_{-i}) = (v_{i,a})_1^2 + (g^{(i)})_1^2 = 1 - p_i(a, \mathbf{a}_{-i}) + C$ where
$C = (g^{(i)})_1^2$ is a constant chosen so all costs are non-negative; we may take $g^{(i)} = 0$ so
$C = 0$.

Set $\Sigma = \mathrm{diag}(1, \infty, \ldots, \infty)$ so only $d_1$ contributes to the Mahalanobis
norm; set $\beta_k = 0$ for all $k$. With this construction,
\[
\BF_i(a \mid \mathbf{a}_{-i}) \;=\; 1 - p_i(a, \mathbf{a}_{-i}).
\]
The augmented-game payoff is $u_i(a, \mathbf{a}_{-i}) = -\BF_i(a \mid \mathbf{a}_{-i}) =
p_i(a, \mathbf{a}_{-i}) - 1$, which is $p_i(a, \mathbf{a}_{-i})$ shifted by a constant. Constant
shifts preserve the best-response structure: agent $i$'s best response in $\Gamma^+$ is the same set
of actions as agent $i$'s best response in $G$, so the $\varepsilon$-Nash equilibria of $G$ are in
one-to-one correspondence with the $\varepsilon$-BGE of $\Gamma$.

The reduction is polynomial-time: each vertex coordinate is computable in time polynomial in the
encoding of $p_i$, and there are $\sum_i |A_i|$ such vertices. The non-negativity of path cost is
preserved because $p_i$ is normalized to $[0,1]$. The construction uses only the diagonal-and-degenerate
covariance and zero boundary penalties, so the manifold-specific structure plays no role in the
hardness. $\square$
\end{proof}

The gadget is tight in the following senses. First, the reduction uses only one active dimension
($d_1$), so \textsc{PPAD}-hardness is preserved even on the most degenerate manifold. Second, the
non-negativity of Mahalanobis distance is respected because the bimatrix payoffs are normalized to
$[0,1]$ before encoding; the Mahalanobis norm becomes $\sqrt{1-p_i(\cdot)} \in [0,1]$, and squaring
gives the additive cost $1 - p_i(\cdot)$. Third, the bimatrix payoff structure is recovered exactly,
not approximately, so an $\varepsilon$-BGE of $\Gamma$ yields an $\varepsilon$-Nash of $G$ for the
same $\varepsilon$.

\begin{remark}[What this result does and does not establish]
\label{rem:scope-of-hardness}
Theorem~\ref{thm:bge-ppad} establishes \textsc{PPAD}-completeness on the \emph{single-action}
manifold-game class. The unrestricted case, in which strategies are arbitrary simple paths through
the decision complex and the set $S_i$ of paths can be exponentially large in the complex size, is a
strictly harder problem: the augmented game is no longer of polynomial size, and the natural
complexity class is probably above \textsc{PPAD}. We do not pursue this case here; the
single-action result already establishes that the manifold framework gains no complexity advantage
over standard Nash in the worst case, which is the point Theorem~\ref{thm:bge-ppad} needs to make.
\end{remark}

\begin{remark}[Why agents cannot just compute BGE]
Theorem~\ref{thm:bge-ppad} formalizes the puzzle that motivates the rest of the paper. Even with the
9-dimensional metric machinery in place, the equilibrium concept inherits the worst-case complexity
of Nash. No actual agent---human, animal, or artificial---can be solving the unrestricted BGE
problem. Yet experimentally validated behavioral equilibria exist and are routinely observed. The
remaining sections explain why.
\end{remark}


%==============================================================================
\section{Heuristic Truncation and Tractability}
\label{sec:truncation}
%==============================================================================

\subsection{Truncation, active sets, and admissibility}

We formalize the truncation operation already implicit in the structural-fuzzing procedure
of~\cite{bond2026tcss}.

\begin{definition}[Active set and truncation]
\label{def:truncation}
For an active set $T \subseteq \{1, \ldots, 9\}$ of dimension indices, the \emph{truncated decision
complex} $\E|_T$ has the same vertex and edge sets as $\E$ but with metric
\[
w_T(v_i, v_j) \;=\; \sqrt{\sum_{k \in T} (\Delta a_k)^2 / \sigma_k^2}
                 \;+\; \sum_{k \in \mathcal{B}(v_i, v_j) \cap T} \beta_k.
\]
Movement on dimensions $k \notin T$ contributes zero cost. The \emph{truncation dimension} is
$k = |T|$.
\end{definition}

In~\cite{bond2026tcss}, the empirical active set is $T = \{6, 7, 9\}$ with $k = 3$.

\begin{definition}[Admissible heuristic truncation]
\label{def:admissible}
A truncation $T$ is \emph{admissible} for the game $\Gamma$ if for every agent $i$ and every node $v$
on agent $i$'s decision complex, the truncated cost-to-go from $v$ to the goal region under $w_T$
underestimates the cost-to-go under the full metric $w$:
\begin{equation}
\label{eq:admissible}
h_T(v) \;\leq\; h^*(v),
\end{equation}
where $h_T(v) = \min_\gamma \sum w_T(v_i, v_{i+1})$ over paths $\gamma$ from $v$ to the goal under
the truncated metric, and $h^*(v)$ is the corresponding minimum under the full metric.
\end{definition}

The terminology is borrowed from the $A^*$ literature \cite{hart1968}: an \emph{admissible
heuristic} is one that never overestimates the true remaining cost. The Heuristic Truncation Theorem
of~\cite[Theorem~7]{bond2026geometric} establishes that the moral-boundary heuristic $h_M$ is
admissible on the full manifold whenever $\beta_k$ underestimates the true cost of crossing
boundary $k$. We use the same notion here, applied to dimension truncation rather than boundary
penalty.

\begin{lemma}[Diagonal truncation preserves admissibility]
\label{lem:diagonal}
If $\Sigma$ is diagonal, then any subset $T \subseteq \{1, \ldots, 9\}$ defines an admissible
truncation in the sense of Definition~\ref{def:admissible}.
\end{lemma}

\begin{proof}
Under a diagonal metric, $w(v_i, v_j)^2 = \sum_{k=1}^{9} (\Delta a_k)^2/\sigma_k^2 + \beta\text{-terms}$.
Restricting the sum to $k \in T$ removes non-negative summands and removes (possibly non-negative)
boundary contributions for $k \notin T$. Thus $w_T(v_i, v_j) \leq w(v_i, v_j)$ for every edge.
Summing along any path and minimizing yields $h_T(v) \leq h^*(v)$. $\square$
\end{proof}

Lemma~\ref{lem:diagonal} is the structural reason the geometric framework supports tractability
arguments: under the diagonality assumption, \emph{every} subset of dimensions is an admissible
truncation. The empirical choice of which subset is made on a different criterion---predictive
accuracy on aggregate behavioral targets, established by structural fuzzing in~\cite{bond2026tcss}.

\subsection{The truncated game and the quotient complex}

\begin{definition}[Truncated game and truncated complex sizes]
\label{def:truncated-game}
Given a manifold game $\Gamma = (N, \{\E_i\}, \{G_i\})$ and an active set $T \subseteq \{1,\ldots,9\}$,
the \emph{truncated game} $\Gamma_T$ has the same agent set $N$ and goal regions $\{G_i\}$, but each
decision complex $\E_i$ is replaced by its truncated counterpart $\E_i|_T$ (Definition~\ref{def:truncation}).
Define the equivalence $v \sim_T v' \iff v_k = v'_k \text{ for all } k \in T$, and let
$\E_i|_T / \sim_T$ denote the resulting \emph{quotient complex}, whose vertex set is the set of
$\sim_T$-equivalence classes and whose edges are the projections of the original edges to the quotient.
Let $V_i^T$ and $E_i^T$ denote the number of vertices and edges of $\E_i|_T / \sim_T$, and set
$V^T = \max_i V_i^T$ and $E^T = \max_i E_i^T$.
\end{definition}

The quotient construction is forced by the truncated metric: under $w_T$, two vertices that agree on
$T$ have identical truncated cost-to-go from every other vertex, so the shortest-path problem on
$\E_i|_T$ is well-defined on the quotient. The number of vertices of the quotient satisfies $V_i^T
\leq V_i$, with equality only if no two vertices coincide on $T$; if the vertex set of $\E_i$ has
product structure across dimensions, $V_i / V_i^T$ can be exponential in $9 - k$. This contraction
of the search graph is the structural source of the truncation's runtime savings, as
Remark~\ref{rem:savings} makes precise below.

Computing a BGE of $\Gamma_T$ reduces to a Nash problem on the augmented truncated game
$\Gamma_T^+$, by the same construction as Theorem~\ref{thm:equivalence}, applied vertex-wise on the
quotient complex.

\subsection{Main result}

We can now state the tractability theorem.

\begin{theorem}[Tractable BGE via admissible heuristic truncation]
\label{thm:tractable-bge}
Let $\Gamma$ be a finite manifold game with $n$ agents and diagonal covariance $\Sigma$. Let
$T \subseteq \{1,\ldots,9\}$ be an active set of dimension $k = |T|$, inducing the truncated game
$\Gamma_T$ and quotient complex sizes $(V^T, E^T)$ as in Definition~\ref{def:truncated-game}.
Suppose:
\begin{enumerate}
\item \textbf{(Contraction).} The truncated game $\Gamma_T$ satisfies the contraction condition of
Lemma~\ref{lem:contraction} with coupling ratio $L < 1$.
\item \textbf{(Heuristic cost).} An admissible heuristic $h_T$ on the truncated complex is
computable in time $t_h = t_h(k)$ per query.
\end{enumerate}
Then for any $\varepsilon > 0$, an $\varepsilon$-approximate BGE of $\Gamma_T$ is computable by
iterated best-response in time
\begin{equation}
\label{eq:complexity}
T_{\mathrm{BGE}}
\;=\;
O\!\left(\; n \cdot \frac{\log(1/\varepsilon)}{\log(1/L)} \cdot \bigl(E^T + V^T \log V^T + t_h(k)\,V^T\bigr) \;\right).
\end{equation}
\end{theorem}

\begin{proof}
\textbf{Iteration count.} By Lemma~\ref{lem:contraction} the joint best-response mapping on
$\Gamma_T^+$ is an $L$-contraction; by Banach's fixed-point theorem the iteration converges
geometrically to the unique fixed point. To achieve precision $\varepsilon$ in the strategy-profile
sup-norm requires $N_{\mathrm{iter}} = \lceil \log(1/\varepsilon)/\log(1/L) \rceil$ iterations.

\textbf{Per-agent best response.} In each iteration, agent $i$'s best response is the shortest
truncated path in $\E_i|_T$ from $v_0^{(i)}$ to $G_i$, equivalently the shortest path on the quotient
complex $\E_i|_T/\sim_T$. By Lemma~\ref{lem:diagonal}, $h_T$ is an admissible heuristic for the
truncated metric on the quotient. The standard analysis of $A^*$ with an admissible
heuristic~\cite{hart1968} establishes that each vertex of the search graph is expanded at most once;
the running time is therefore
\[
O\!\left(E_i^T + V_i^T \log V_i^T + t_h(k)\cdot V_i^T\right),
\]
where the first two terms are the priority-queue bound and the third counts heuristic evaluations.

\textbf{Total runtime.} Summing across $n$ agents per iteration and across $N_{\mathrm{iter}}$
iterations gives~\eqref{eq:complexity}. $\square$
\end{proof}

\begin{corollary}[Polynomial-time approximate BGE at fixed $k$]
\label{cor:polytime}
For any fixed $k$ and any uniform bound $t_h(k) \leq \mathrm{poly}(k)$, the truncated BGE of a finite
manifold game with contractive coupling can be $\varepsilon$-approximated in time polynomial in $n$,
$V^T$, $E^T$, and $\log(1/\varepsilon)$. The dependence on $L$ enters through $1/\log(1/L)$.
\end{corollary}

\begin{remark}[Where the truncation savings come from]
\label{rem:savings}
The previous draft of this paper attributed an exponential factor $2^{O(k)}$ to the worst-case
branching of $A^*$. That bound is incorrect: $A^*$ with an \emph{admissible} heuristic expands each
vertex at most once~\cite{hart1968}, so the per-call runtime is linear in $V^T$ up to the
priority-queue and heuristic-evaluation overheads. The truncation's runtime benefit is therefore not
an $A^*$-branching argument. It comes from three structurally distinct sources:
\begin{enumerate}[label=(\alph*)]
\item \textbf{Quotient graph contraction.} States that differ only on inactive dimensions are identified
under $\sim_T$, so the search graph shrinks from $(V, E)$ to $(V^T, E^T)$ where the ratio $V/V^T$ can
be exponential in $9-k$ for product-structured decision complexes. This is the dominant savings.
\item \textbf{Heuristic-evaluation cost.} The heuristic computation $t_h(k)$ depends on $k$ rather than
$9$, so the per-vertex overhead scales with the active dimension count.
\item \textbf{Coupling ratio.} The contraction ratio $L$ in the truncated game can be strictly
smaller than the ratio in the untruncated game, because coupling channels passing through inactive
dimensions are removed. We do not prove a general inequality here; the empirical anchor in
Section~\ref{sec:empirical} reports calculated values for the calibrated regime as indicative
diagnostics.
\end{enumerate}
\end{remark}

\begin{remark}[Bypass on a restricted class, not a defeat of \textsc{PPAD}]
\label{rem:bypass}
Theorem~\ref{thm:tractable-bge} does \emph{not} defeat \textsc{PPAD}-hardness. The hardness of
Theorem~\ref{thm:bge-ppad} holds for the unrestricted single-action manifold-game class. The
tractability result holds for the strictly smaller class of \emph{contractive} manifold games under
an \emph{admissible truncation} that meaningfully shrinks the quotient complex. That contraction
conditions yield tractability for fixed-point computation is a known result in algorithmic game theory
(cf.\ supermodularity arguments~\cite{milgrom1990} and potential-game algorithms~\cite{monderer1996}).
What is new in Theorem~\ref{thm:tractable-bge} is not the use of contraction \emph{per se} but
the identification of \emph{which class} satisfies it: the empirically validated behavioral
equilibrium class identified by the structural-fuzzing calibration of~\cite{bond2026tcss}. The
restricted class is what empirical behavior actually populates; the contraction condition is what
makes that empirical region tractable; the conjunction is what locates behavioral game theory inside
the polynomial-time fragment of equilibrium computation. The novelty is the joining, not either side
alone.
\end{remark}

\subsection{Approximation error of the truncation}

The constant $L$ in Theorem~\ref{thm:tractable-bge} controls the convergence speed of best-response
iteration; the truncated complex sizes $(V^T, E^T)$ and the heuristic-evaluation cost $t_h(k)$
control the per-iteration cost. A separate question is the \emph{approximation error} of the
truncated BGE $\gamma_T^*$ relative to the BGE $\gamma^*$ of the full (untruncated) game $\Gamma$.
The next proposition bounds it explicitly.

\begin{proposition}[Approximation error of truncated BGE]
\label{prop:approx}
Let $\gamma^*$ denote the BGE of $\Gamma$ on the full manifold and $\gamma_T^*$ the BGE of
$\Gamma_T$. Then for every agent $i$,
\[
\bigl| \BF_i(\gamma_T^*) - \BF_i(\gamma^*) \bigr|
\;\leq\; \max_{k \notin T} \sigma_k^{-2} \cdot \|\Delta\mathbf{a}^{(i)}\|_\infty^2 \cdot |\gamma_i^*|,
\]
where $|\gamma_i^*|$ denotes the number of edges along agent $i$'s equilibrium path.
\end{proposition}

The proof follows from the diagonal-metric bound on the per-edge cost contribution of the inactive
dimensions. Proposition~\ref{prop:approx} formalizes a small-perturbation guarantee: if the
inactive-dimension variances $\sigma_k^2$ for $k \notin T$ are large, the truncation error is small.
When $\sigma_k^2 = \infty$ for all $k \notin T$, the truncation is exact: $\gamma_T^* = \gamma^*$.

This last observation matters for the empirical anchor. The TCSS calibration sets
$\sigma_k^2 = \infty$ on the six inactive dimensions, so the truncated BGE coincides with the BGE on
the full manifold under the calibrated metric. There is no approximation error within the calibrated
regime---only the question of whether the regime itself is the right one.


%==============================================================================
\section{Empirical Anchor: the TCSS Calibration}
\label{sec:empirical}
%==============================================================================

We verify that the empirically validated calibration of~\cite{bond2026tcss} instantiates the
hypotheses of Theorem~\ref{thm:tractable-bge}.

\subsection{The active set and the truncation}

\cite{bond2026tcss} apply structural fuzzing---an exhaustive search over all $\sum_{k=1}^{5}
\binom{9}{k} = 381$ candidate subsets of dimensions---to minimize aggregate prediction MAE on nine
game-theoretic calibration targets. The selected active set is $T^\star = \{6, 7, 9\}$
(Social Impact, Virtue/Identity, Epistemic Status) with calibrated variances $\sigma_6^2 = 78.26$,
$\sigma_7^2 = 32.28$, $\sigma_9^2 = 0.01262$. All other dimensions are inactive
($\sigma_k^2 = \infty$).

By Lemma~\ref{lem:diagonal}, $T^\star$ is an admissible truncation. By the active-set structure,
Proposition~\ref{prop:approx} bounds the truncation error at zero within the calibration regime: the
truncated BGE \emph{is} the BGE on the full manifold under the calibrated metric.

\subsection{Contraction analysis: scope and the single-shot calibration regime}

Theorem~\ref{thm:tractable-bge} requires the contraction condition (Lemma~\ref{lem:contraction}) on
the truncated game. Appendix~\ref{app:lipschitz} carries out the rigorous derivation of the
Lipschitz bound on the joint best-response mapping using the Boltzmann--Gibbs operator-norm
estimate: for a softmax best-response with temperature $T$ and cost-coupling Lipschitz constant
$K_i$ (Assumption~\ref{ass:cost-coupling}), the joint BR satisfies
\[
\bigl\| \mathrm{BR}(\sigma) - \mathrm{BR}(\sigma') \bigr\|_{\mathrm{TV}}
\;\leq\; L\,\bigl\| \sigma - \sigma' \bigr\|_{\mathrm{TV}},
\qquad L \;=\; \max_i \frac{K_i}{2T}.
\]
The contraction condition $L < 1$ thus reduces to $K_i < 2T$ for every agent~$i$
(Proposition~\ref{prop:br-lip}).

The TCSS empirical calibration is, however, \emph{single-shot} rather than iterated. It uses a fixed
empirical responder rejection logistic
$P_{\text{reject}}(x) = 1/(1 + \exp((x - x_{\text{thresh}})/k))$ with $x_{\text{thresh}} = 0.18$,
$k = 0.15$, estimated from the Fraser--Nettle ultimatum data~\cite{bond2026tcss}, and the proposer's
strategy is the one-shot $A^*$-optimal response against this fixed function. There is no iterated
best-response in the calibration; the contraction condition is trivially saturated (the iteration
converges in one step), and the runtime bound~\eqref{eq:complexity} reduces to a single $A^*$ call
per agent. Appendix~\ref{app:lipschitz}.5 develops this distinction carefully.

The contraction bound from Appendix~\ref{app:lipschitz} is therefore best read as a forward-going
prediction about iterated extensions of the framework---trust games with reputation updating,
public-goods games with conditional cooperation, network games with reinforcement-learning
agents---in which the rigorous Lipschitz machinery applies with the same calibrated parameters. The
explicit bound under the calibrated metric is
$L \leq 1/(2T) \leq 1$ (Appendix~\ref{app:lipschitz}.3), which saturates the contraction boundary
in total-variation norm; a strict contraction in the calibrated regime is established under a
dominance-of-cost-over-goal assumption (Appendix~\ref{app:lipschitz}.4) and is an open question in
the unrestricted case.

\subsection{The complexity bound under the calibrated truncation}

In the single-shot TCSS calibration regime, the BGE reduces to a per-agent $A^*$ call on the
truncated complex with calibrated metric. The runtime bound~\eqref{eq:complexity} simplifies
because the iteration count is one ($L \to 0$ for non-iterated games), and the only contributions
are the truncated complex sizes and the heuristic evaluation cost:
\[
T_{\mathrm{BGE}}^{\mathrm{TCSS}}
\;=\;
O\!\left(\; n \cdot \bigl(E^T + V^T \log V^T + t_h(3)\,V^T\bigr) \;\right).
\]
By the explicit counts in Appendix~\ref{app:vt-et} (Table~\ref{tab:per-target-vt}), $V^T \leq 13$
and $E^T \leq 22$ across all sixteen calibration targets, and the diagonal-Mahalanobis heuristic
has $t_h(3) = O(1)$. The per-agent cost is therefore a small constant, dominated by $\approx 70$
elementary operations per $A^*$ call. The PPAD-hard worst case of Theorem~\ref{thm:bge-ppad}
becomes a polynomial-time computation (in fact constant-time per agent) on this restricted
behavioral subclass; we emphasize ``bypassed,'' not ``defeated.''

For genuine iterated extensions---trust games, repeated public-goods games, network games---the
full bound~\eqref{eq:complexity} applies with the iteration count $\log(1/\varepsilon)/\log(1/L)$
controlled by the Lipschitz analysis of Appendix~\ref{app:lipschitz}.

\subsection{The high-stakes falsification: empirical adequacy, not formal admissibility}

The framework's stake-scaling-invariance corollary, when probed on Andersen et al.'s
\cite{andersen2011} high-stakes Indonesia ultimatum data ($n{=}458$, stakes Rs~20 to Rs~20{,}000),
is rejected at $p < 10^{-3}$ on all three pre-specified tests (Kruskal--Wallis on offers, $\chi^2$
on rejections, logit LRT on rejection conditional on offer share)~\cite{bond2026tcss}. We are
careful about what this implies for the present framework.

\textbf{Formal admissibility is unaffected.} By Lemma~\ref{lem:diagonal}, $T^\star = \{6, 7, 9\}$
remains an admissible truncation in the $A^*$ sense for any positive stake: $h_{T^\star} \leq h^*$
holds identically because removing non-negative summands from a Mahalanobis squared sum cannot
increase it. Admissibility is a property of the truncation itself, not of the game's stake structure.

\textbf{Empirical adequacy fails.} What the Andersen result establishes is that the
\emph{approximation-error bound of Proposition~\ref{prop:approx}} becomes too loose to be predictive.
Concretely, the bound contains $\max_{k \notin T^\star} \sigma_k^{-2}$, which under the calibrated
metric is zero (the inactive variances are infinite). The bound therefore says the truncated BGE
\emph{equals} the BGE on the full manifold under the calibrated metric. The Andersen rejection shows
that the calibrated metric itself is wrong at high stakes: a recalibration with finite $\sigma_1^2$
is required. After recalibration, the new metric induces a non-zero Proposition~\ref{prop:approx}
bound that quantifies the residual truncation error at the new active set, and the boundary at
which this bound becomes acceptable defines the regime of empirical adequacy.

\textbf{What this distinguishes.} Three notions must be kept separate:
\begin{enumerate}[label=(\roman*)]
\item \emph{Formal admissibility of the truncation:} $h_T \leq h^*$ holds always (Lemma~\ref{lem:diagonal}).
\item \emph{Approximation tightness:} the bound of Proposition~\ref{prop:approx} is non-vacuous only when
the inactive variances are finite; it is vacuously tight (equal-to-zero) when $\sigma_k^2 = \infty$
for $k \notin T$.
\item \emph{Empirical adequacy of the calibrated metric:} the calibrated $\Sigma$ may itself be wrong
in regimes outside the calibration sample, in which case (ii) is the wrong measurement to make.
\end{enumerate}
The Andersen rejection lies in (iii), not (i) or (ii). The framework's response is to recalibrate
$\Sigma$ on the new regime, not to claim the truncation has lost its $A^*$-admissibility property.
We treat the construction of the recalibrated metric---finding the smallest active set $T'$ that
contains $T^\star$ and restores empirical fit on high-stakes data---as the principal target of
follow-up work.


%==============================================================================
\section{Bounded Rationality, Reframed}
\label{sec:reframe}
%==============================================================================

Theorem~\ref{thm:tractable-bge} suggests a reframing of the bounded-rationality literature. We sketch
the implications and contrast the framework with several proximate alternatives.

\subsection{Simon's satisficing}

\cite{simon1955} characterized bounded rationality as satisficing: agents accept any alternative
above an aspiration level rather than searching for the optimum. Satisficing is a stopping rule;
heuristic truncation is a search-space restriction. The two are complementary: satisficing terminates
$A^*$ early on the truncated complex; truncation determines which complex $A^*$ runs on. Both are
necessary for tractable behavior under PPAD-hardness, but they answer different questions. Simon's
question was \emph{when} to stop; ours is \emph{where} to search.

\subsection{Heuristics-and-biases}

The heuristics-and-biases programme \cite{tversky1974,kahneman1979} catalogued systematic deviations
from expected-utility maximization. The standard interpretation reads these as cognitive failures.
The geometric framework reads them as admissible heuristics on the full manifold whose deviations
from rationality appear only in the scalar projection~\cite[Theorem~10]{bond2026geometric}; the
present paper adds the complexity-theoretic reason: deviation is the cost of tractability. An agent
that committed to global optimization on the full manifold would face the PPAD barrier and would
exhibit no behavior at all in finite time.

\subsection{Quantal-response equilibrium}

\cite{mckelvey1995} introduced the quantal-response equilibrium (QRE) as a smoothed Nash concept in
which agents play noisy best responses, with the noise scaled by a temperature parameter. QRE and
truncated BGE share a stochastic-choice softmax structure. The conceptual difference is that QRE
treats the noise as a free parameter to be fit, whereas the geometric framework derives the choice
distribution from a cost-dependent temperature
$T(\Delta) = \max(T_{\text{floor}}, T_{\text{base}} \Delta^{T_\alpha})$ defined on the
manifold~\cite{bond2026tcss}. The empirical fitted values
$(T_{\text{base}}, T_\alpha, T_{\text{floor}}) = (0.24, 2.13, 0.5)$ are calibration constants of the
geometric model, not free QRE parameters.

\subsection{Resource rationality}

\cite{lieder2020} formalized bounded rationality as resource-constrained optimization: agents
maximize expected utility minus computational cost. The framework is compatible with ours but
operates at a different level. Resource rationality predicts \emph{how much} computation an agent
should perform; the geometric framework specifies \emph{which structure} the computation is performed
on, and the present paper shows why the latter choice has complexity-theoretic teeth: without
heuristic truncation, the agent has no tractable equilibrium to converge to, regardless of how much
compute it is willing to spend.


%==============================================================================
\section{Limitations and Open Questions}
\label{sec:limitations}
%==============================================================================

The result is narrow in several respects.

\begin{enumerate}

\item \textbf{Diagonal covariance.} Lemma~\ref{lem:diagonal} relies on the diagonality of $\Sigma$.
Off-diagonal interactions would couple dimensions and could violate admissibility under naive
truncation. The geometric framework's general formulation allows non-diagonal $\Sigma$, but the
tractability theorem covers only the diagonal restriction. Extending the result to sparse off-diagonal
structure (where the inactive--inactive and active--inactive blocks vanish) is a tractable
follow-up; the active--active block could induce additional contraction or anti-contraction effects
that require separate analysis.

\item \textbf{Selecting the active set.} The empirical TCSS active set $\{6, 7, 9\}$ is selected by
structural fuzzing on aggregate behavioral data, not derived from first principles. The present paper
takes the active set as input. Whether there is a principled procedure for predicting which subset
an arbitrary game will admit as its admissible truncation remains open. Candidates include
adversarial-encoding probes, factor-analytic decompositions of multi-task individual-level choice
data, and external-rater encoding validations.

\item \textbf{Heuristic accuracy and quotient size.} Theorem~\ref{thm:tractable-bge} parametrizes
runtime by the truncated complex sizes $(V^T, E^T)$ and the heuristic-evaluation cost $t_h(k)$. Tighter
heuristics reduce the constant absorbed into $t_h$, and product-structured decision complexes give
$V^T \ll V$ under the quotient construction. The empirical TCSS targets exhibit small $V$ to begin
with (action sets of size $\leq 11$) so the quotient savings are bounded; targets with richer action
sets would benefit more from the quotient construction. A refined complexity bound that depends on
the heuristic gap $h^* - h_T$ is open.

\item \textbf{Mixed BGE.} Theorem~\ref{thm:tractable-bge} treats convergence to pure BGE under
contraction. Mixed BGE with pure-strategy cycling requires the full Nash machinery and inherits the
\textsc{PPAD} complexity of mixed equilibria. The contraction condition rules out cycles; understanding
the trade-off between admissible truncation and mixed-equilibrium support is open.

\item \textbf{Multi-agent scaling.} The number of agents $n$ enters linearly in \eqref{eq:complexity}
through best-response cost per iteration. Many-agent settings with sparse coupling structure
(networks, marketplaces) admit further reductions, but a formal extension of the contraction lemma
to network-structured games is not pursued here.

\item \textbf{Empirical adequacy of the calibrated metric.} The TCSS test on the
Andersen et al.~\cite{andersen2011} data identifies one regime in which the calibrated metric is
empirically inadequate (see Section~\ref{sec:empirical}.4 for the distinction between formal
admissibility and empirical adequacy). A systematic map of regimes---across stake levels, cultural
contexts, institutional settings, and game classes---would convert the truncation procedure from
descriptive to predictive.

\end{enumerate}


%==============================================================================
\section{Conclusion}
\label{sec:conclusion}
%==============================================================================

Computing a Nash equilibrium is \textsc{PPAD}-complete. By the BGE--Nash equivalence, computing a
Bond Geodesic Equilibrium of a single-action manifold game is also \textsc{PPAD}-complete
(Theorem~\ref{thm:bge-ppad}). Yet aggregate behavioral equilibria are routinely observed in human
and artificial agents. This paper closes the gap not by defeating \textsc{PPAD} but by locating the
behavioral region inside the polynomial-time fragment: under admissible heuristic truncation to a
small active set and a contraction condition on the truncated coupling, the BGE is computable by
iterated best-response in time polynomial in the agents, the truncated complex sizes, and
$\log(1/\varepsilon)$ (Theorem~\ref{thm:tractable-bge}).

The reframing is constructive: bounded rationality is the computational price of tractable
equilibrium. The geometric framework's interpretable dimensions are not merely descriptive labels for
behavioral aggregates but candidate heuristic coordinates over which admissible truncation is
performed. The empirically validated 3-dimensional active set
$\{d_6, d_7, d_9\}$~\cite{bond2026tcss} instantiates the theorem's hypotheses and bounds its
complexity at a constant factor in the truncation dimension.

Future work points in several directions. Off-diagonal covariance structure, principled selection of
the active set, neural and algorithmic implementations of admissible truncation in agent
architectures, and the systematic mapping of regimes in which calibrated truncations remain
empirically adequate are all open. The result demonstrated here is structural: the \textsc{PPAD}
barrier does not preclude tractable equilibrium on the restricted class behavioral data actually
populate, provided the equilibrium concept is the right one and the search space is truncated to
coordinates the agent can actually navigate.

%==============================================================================
\appendix
\section{Quotient Complex Sizes for the TCSS Calibration Targets}
\label{app:vt-et}
%==============================================================================

Theorem~\ref{thm:tractable-bge} parametrizes the per-iteration runtime by the truncated complex
sizes $V^T$ and $E^T$ (Definition~\ref{def:truncated-game}) and the heuristic-evaluation cost
$t_h(k)$. This appendix computes these quantities explicitly for each of the sixteen targets in the
TCSS calibration~\cite{bond2026tcss} under the active set $T^\star = \{6, 7, 9\}$.

\subsection*{A.1 Per-target action sets}

Each calibration target is a single-decision problem with a discrete candidate action set
$A_i$. The decision complex $\E_i$ has the canonical single-action structure: one source vertex
$v_0$, one vertex $v_{i,a}$ per candidate action $a \in A_i$, and one terminal vertex $g$,
connected by edges $v_0 \to v_{i,a} \to g$.

The candidate action sets used in the TCSS calibration are:
\begin{itemize}
\item \textbf{Ultimatum game} (proposer mean offer, proposer modal offer, G\"uth 1982 replication).
Offer fraction $x$ on the grid $\{0, 0.05, 0.10, \ldots, 0.50\}$, so $|A_i| = 11$.
\item \textbf{Dictator game.} Same offer grid as ultimatum: $|A_i| = 11$.
\item \textbf{Ultimatum responder.} Binary choice (accept or reject) at each candidate offer
threshold. Aggregated to the population-level minimum acceptable offer (MAO) by integrating the
rejection logistic, but the underlying single-decision structure has $|A_i| = 2$.
\item \textbf{Public-goods game} (rounds 1, 3, 5, 8, 10). Contribution fraction $f$ on the grid
$\{0, 0.10, 0.20, \ldots, 1.00\}$, so $|A_i| = 11$ per round.
\item \textbf{Prospect-theory items} (P1, P3, P7, P11, P16, P17). Binary choice between two
lotteries: $|A_i| = 2$.
\end{itemize}

\subsection*{A.2 Per-target counts}

For each target, $V_i = |A_i| + 2$ (source and goal plus action vertices) and $E_i = 2|A_i|$.
Under $T^\star = \{6, 7, 9\}$, the equivalence $v \sim_T v'$ identifies vertices whose
coordinates on $T^\star$ coincide. From the encoding rules of~\cite[\S IV]{bond2026tcss}
(reproduced verbatim in Section~\ref{sec:empirical} above), distinct candidate actions give
distinct values on the $T^\star$ coordinates: the ultimatum/dictator encoding has $d_6 = d_7 = x$,
each $x$ unique; the public-goods encoding has $d_6 = f \cdot \delta(r)$ and $d_7 = f$ which
jointly identify $f$ uniquely; the prospect-theory encoding produces two lotteries with distinct
$(d_6, d_7, d_9)$ values. The source and goal have $d_6 = d_7 = 0$ and so collapse into one
$\sim_T$-class each. The truncated counts therefore equal the untruncated counts on the action
vertices; we report both for clarity in Table~\ref{tab:per-target-vt}.

\begin{table}[h]
\centering
\renewcommand{\arraystretch}{1.1}
\caption{Decision-complex sizes per calibration target.
$|A_i|$ is the action-set size, $V_i$ and $E_i$ are the untruncated vertex/edge counts,
$V_i^T$ and $E_i^T$ are the quotient sizes under $T^\star = \{6, 7, 9\}$.}
\label{tab:per-target-vt}
\begin{tabular}{l c c c c c}
\toprule
\textbf{Target} & $|A_i|$ & $V_i$ & $E_i$ & $V_i^T$ & $E_i^T$ \\
\midrule
Ultimatum mean offer       & 11 & 13 & 22 & 13 & 22 \\
Ultimatum modal offer      & 11 & 13 & 22 & 13 & 22 \\
Dictator giving            & 11 & 13 & 22 & 13 & 22 \\
Ultimatum responder MAO    &  2 &  4 &  4 &  4 &  4 \\
PG round 1                 & 11 & 13 & 22 & 13 & 22 \\
PG round 3                 & 11 & 13 & 22 & 13 & 22 \\
PG round 5                 & 11 & 13 & 22 & 13 & 22 \\
PG round 8                 & 11 & 13 & 22 & 13 & 22 \\
PG round 10                & 11 & 13 & 22 & 13 & 22 \\
PT P1 (Allais)             &  2 &  4 &  4 &  4 &  4 \\
PT P3 (Strong certainty)   &  2 &  4 &  4 &  4 &  4 \\
PT P7 (Reflection)         &  2 &  4 &  4 &  4 &  4 \\
PT P11 (Isolation)         &  2 &  4 &  4 &  4 &  4 \\
PT P16 (Small-prob gain)   &  2 &  4 &  4 &  4 &  4 \\
PT P17 (Small-prob loss)   &  2 &  4 &  4 &  4 &  4 \\
G\"uth 1982 replication    & 11 & 13 & 22 & 13 & 22 \\
\midrule
\textbf{Maximum}           & 11 & 13 & 22 & 13 & 22 \\
\bottomrule
\end{tabular}
\end{table}

\subsection*{A.3 What these numbers do, and do not, demonstrate}

Three observations matter for interpreting Table~\ref{tab:per-target-vt}.

\textbf{(i) The quotient construction is vacuous on these targets.}
For every calibration target, $V_i^T = V_i$ and $E_i^T = E_i$. The reason is that the TCSS
encoding rules make $d_6$ and $d_7$ (and, in the prospect-theory items, the lottery-dependent
$d_9$) distinct functions of the chosen action, so no two action vertices coincide on the
$T^\star$ coordinates. The hypothetical exponential savings from quotient contraction sketched in
Remark~\ref{rem:savings}(a) do not materialize here. They would materialize on richer complexes
in which states differing only on inactive dimensions exist as separate vertices (for example, a
complex enumerating combinations of offer amount, side-payment amount, and timing; the side-payment
and timing dimensions could be inactive while offer remains active, producing $V/V^T$ exponential
in the number of inactive choice axes).

\textbf{(ii) The empirical savings come from the heuristic-evaluation cost and the contraction
ratio.} With $V^T \leq 13$ and $E^T \leq 22$ across all targets, the per-iteration cost of
iterated best-response under Theorem~\ref{thm:tractable-bge} is
\[
E^T + V^T \log V^T + t_h(3) \cdot V^T
\;\;\leq\;\;
22 + 13 \log 13 + 13\, t_h(3)
\;\;\approx\;\;
70 + 13\, t_h(3),
\]
which is dominated by the heuristic-evaluation term. Under the diagonal-Mahalanobis heuristic
$h_T(v) = \sqrt{\sum_{k \in T} (\Delta a_k)^2/\sigma_k^2}$ from~\cite{bond2026tcss}, $t_h(3) = O(1)$
(three multiplications, two additions, one square root). The per-iteration cost is therefore a
small constant in the calibrated regime, and the total runtime of~\eqref{eq:complexity} reduces to
\[
T_{\mathrm{BGE}}^{\mathrm{TCSS}}
\;=\;
O\!\left(\; n \cdot \frac{\log(1/\varepsilon)}{\log(1/L)} \;\right).
\]
This is the explicit form Section~\ref{sec:empirical}.3 reports as ``the per-iteration cost is
dominated by a constant.''

\textbf{(iii) The bound is tight for sequential best-response on the calibrated complexes.}
Because $V^T$ is small, the constant-factor overhead of $A^*$ versus Dijkstra is negligible; either
algorithm achieves the per-iteration bound. The contraction ratio $L$ is the only non-trivial
parameter affecting runtime in the calibrated regime, and its indicative value $L \approx 0.25$
(Section~\ref{sec:empirical}.2) gives $\log(1/\varepsilon)/\log(1/L) \approx 0.72 \log(1/\varepsilon)$,
so achieving $\varepsilon = 10^{-3}$ takes approximately $5$ best-response iterations and
$\varepsilon = 10^{-6}$ takes approximately $10$.

\subsection*{A.4 Worked example: ultimatum-game proposer}

To make the construction concrete, we trace it for the ultimatum proposer target. The candidate
action set is $A = \{0, 0.05, 0.10, \ldots, 0.50\}$. For each $x \in A$, the action vertex $v_x$ has
9-d coordinates $(\Delta a_1, \ldots, \Delta a_9) = (-x, 0, -(0.5-x)^2, 0, 0, x, x, 0, 0.80)$
(from~\cite[\S IV.A]{bond2026tcss}). The decision complex is the bipartite-plus-source-and-goal
structure $\{v_0\} \cup \{v_x : x \in A\} \cup \{g\}$ with edges $v_0 \to v_x \to g$ for each
$x \in A$.

Under $T^\star = \{6, 7, 9\}$, the projection retains $(d_6, d_7, d_9) = (x, x, 0.80)$. The eleven
action vertices project to the eleven points $\{(0, 0, 0.80), (0.05, 0.05, 0.80), \ldots,
(0.50, 0.50, 0.80)\}$, all distinct. The source projects to $(0, 0, 0)$ and the goal to its own
coordinates (typically $(0, 0, 0)$ when goal is ``offer made''). Under $\sim_T$, no two action
vertices collapse, so $V_i^T = V_i = 13$ and $E_i^T = E_i = 22$. The path cost from $v_0$ to $g$
through $v_x$ is the Mahalanobis sum
\[
w_T(v_0, v_x) + w_T(v_x, g)
\;=\;
2\sqrt{x^2/\sigma_6^2 + x^2/\sigma_7^2 + (0.80)^2/\sigma_9^2}
\]
plus the rejection-risk term (Section~\ref{sec:empirical}.2). The proposer's best response is the
$x \in A$ minimizing this expression given the responder's current strategy. With $V^T = 13$, the
$A^*$ (or equivalently Dijkstra) call returns the minimum in time linear in $V^T$. The total
runtime of best-response iteration is dominated by $L$ and $\log(1/\varepsilon)$.

\subsection*{A.5 When the appendix's bounds would tighten or loosen}

The bounds in Table~\ref{tab:per-target-vt} are derived from the TCSS-published action grids. They
would tighten if the action grid is coarser (e.g., five offer values rather than eleven) and loosen
if it is finer (e.g., a continuous action space approximated by a 100-point grid would give
$V_i^T = V_i = 102$, still tractable but with a larger constant). For the public-goods game with
$n$ agents and per-round contribution choices, the augmented-game payoff $u_i$ depends on the
profile of all $n$ contributions, so the iterated best-response cost scales as $n \cdot V^T$ per
iteration; this is the dominant scaling for many-agent settings.

%==============================================================================
\section{Lipschitz Analysis of the Best-Response Mapping}
\label{app:lipschitz}
%==============================================================================

This appendix provides the rigorous derivation of the contraction ratio $L$ used in
Theorem~\ref{thm:tractable-bge}. The derivation proceeds in three steps: (B.1) the
Boltzmann--Gibbs operator-norm bound on the softmax map, (B.2) the composition with the
opponent-dependent cost map, (B.3) the resulting Lipschitz constant for the joint best-response.
We then apply the bound to (B.4) the iterated ultimatum game and (B.5) clarify the
relationship between the rigorous bound and the TCSS empirical calibration regime, which is
single-shot rather than iterated.

\subsection*{B.1 The Boltzmann--Gibbs Lipschitz bound}

Let $A$ be a finite action set, $c : A \to \R$ a cost function, and $T > 0$ a temperature. The
softmax distribution is
\begin{equation}
\label{eq:softmax-def}
P_T(a \mid c) \;=\; \frac{\exp(-c(a)/T)}{\sum_{a' \in A} \exp(-c(a')/T)}.
\end{equation}

\begin{lemma}[Boltzmann--Gibbs Lipschitz bound]
\label{lem:softmax-lip}
For any two cost functions $c, c' : A \to \R$,
\begin{equation}
\label{eq:softmax-lip}
\bigl\| P_T(\cdot \mid c) - P_T(\cdot \mid c') \bigr\|_{\mathrm{TV}}
\;\leq\; \frac{1}{2T}\, \bigl\| c - c' \bigr\|_\infty,
\end{equation}
where $\|\cdot\|_{\mathrm{TV}}$ is total variation and $\|\cdot\|_\infty$ is the sup norm on $A$.
\end{lemma}

\begin{proof}
The map $c \mapsto -c/T$ is linear with operator norm $1/T$ from $\ell^\infty(A)$ to itself. The
softmax map $u \mapsto \exp(u)/\sum_{a'} \exp(u(a'))$ from $\ell^\infty(A)$ to the simplex
$\Delta(A)$ has Lipschitz constant $1/2$ in the TV norm: this is a standard result for the
log-sum-exp function~\cite{boyd2004}. Composing gives Lipschitz constant $1/(2T)$. $\square$
\end{proof}

The factor $1/(2T)$ is sharp in this norm pairing. The bound is independent of $|A|$, which matters
for the per-iteration scaling.

\subsection*{B.2 Coupling through opponent-dependent cost}

Now consider a two-agent game with strategy profiles $(\sigma_1, \sigma_2) \in \Delta(A_1) \times
\Delta(A_2)$. Each agent $i$'s cost depends on the opponent's mixed strategy through a function
$c_i(a_i \mid \sigma_{-i})$. The best-response mapping is the softmax
$\mathrm{BR}_i(\sigma_{-i})(a_i) = P_{T}(a_i \mid c_i(\cdot \mid \sigma_{-i}))$.

\begin{assumption}[Cost coupling Lipschitz constant]
\label{ass:cost-coupling}
For each agent $i$, the cost function $\sigma_{-i} \mapsto c_i(\cdot \mid \sigma_{-i})$ is Lipschitz
from $(\Delta(A_{-i}), \|\cdot\|_{\mathrm{TV}})$ to $(\R^{A_i}, \|\cdot\|_\infty)$ with constant
$K_i$:
\begin{equation}
\label{eq:cost-lip}
\sup_{a_i \in A_i} \bigl| c_i(a_i \mid \sigma_{-i}) - c_i(a_i \mid \sigma_{-i}') \bigr|
\;\leq\; K_i \, \bigl\| \sigma_{-i} - \sigma_{-i}' \bigr\|_{\mathrm{TV}}.
\end{equation}
\end{assumption}

\begin{proposition}[Best-response Lipschitz bound]
\label{prop:br-lip}
Under Assumption~\ref{ass:cost-coupling}, the joint best-response mapping
$\mathrm{BR} = (\mathrm{BR}_1, \mathrm{BR}_2)$ on $\Delta(A_1) \times \Delta(A_2)$ satisfies
\begin{equation}
\label{eq:br-lip}
\bigl\| \mathrm{BR}(\sigma) - \mathrm{BR}(\sigma') \bigr\|_{\mathrm{TV}}
\;\leq\; L\, \bigl\| \sigma - \sigma' \bigr\|_{\mathrm{TV}},
\qquad L \;=\; \max_{i \in \{1,2\}} \frac{K_i}{2T}.
\end{equation}
Therefore, if $K_i < 2T$ for both $i$, the BR mapping is a contraction on the strategy-profile
simplex with constant $L < 1$, and by Banach's fixed-point theorem a unique BGE exists and iterated
best-response converges geometrically.
\end{proposition}

\begin{proof}
For agent $i$,
\begin{align*}
\bigl\| \mathrm{BR}_i(\sigma_{-i}) - \mathrm{BR}_i(\sigma_{-i}') \bigr\|_{\mathrm{TV}}
&\;\leq\; \frac{1}{2T}\, \bigl\| c_i(\cdot \mid \sigma_{-i}) - c_i(\cdot \mid \sigma_{-i}') \bigr\|_\infty
& \text{(Lemma~\ref{lem:softmax-lip})} \\
&\;\leq\; \frac{K_i}{2T}\, \bigl\| \sigma_{-i} - \sigma_{-i}' \bigr\|_{\mathrm{TV}}
& \text{(Assumption~\ref{ass:cost-coupling})} \\
&\;\leq\; \frac{K_i}{2T}\, \bigl\| \sigma - \sigma' \bigr\|_{\mathrm{TV}}.
\end{align*}
Taking the maximum over agents gives the joint Lipschitz constant $L$. $\square$
\end{proof}

\subsection*{B.3 The coupling constants under the TCSS encoding}

We compute $K_i$ explicitly for the iterated ultimatum game with the TCSS calibrated metric
under truncation $T^\star = \{6, 7, 9\}$.

\textbf{Proposer's coupling.} The proposer's effective cost at offer $x$, given the responder's
threshold distribution $F_r$ (the mixed strategy over thresholds $\theta \in A_r$), is
\begin{equation}
\label{eq:proposer-cost}
c_p(x \mid F_r) \;=\; \underbrace{\sqrt{x^2 (\sigma_6^{-2} + \sigma_7^{-2}) + 0.80^2 \sigma_9^{-2}}}_{\text{truncated Mahalanobis cost}}
\;-\; (1-x)\, P_{\text{accept}}(x \mid F_r),
\end{equation}
where $P_{\text{accept}}(x \mid F_r) = \mathbb{P}_{\theta \sim F_r}[\theta \leq x] = F_r(x)$ is the
acceptance probability under the responder's threshold CDF. The second term encodes the proposer's
goal-reachability constraint: a rejected offer does not reach the ``offer accepted'' goal, so the
expected payoff $(1-x) P_{\text{accept}}$ enters as a negative cost.

The first term in~\eqref{eq:proposer-cost} is independent of $F_r$. The second satisfies
\[
\bigl| (1-x)\, F_r(x) - (1-x)\, F_r'(x) \bigr|
\;\leq\; (1-x)\, |F_r(x) - F_r'(x)|
\;\leq\; |F_r(x) - F_r'(x)|
\;\leq\; \bigl\| F_r - F_r' \bigr\|_{\mathrm{TV}},
\]
using $|1-x| \leq 1$ on the action grid $x \in [0, 0.5]$ and the fact that $|F_r(x) - F_r'(x)|
\leq \|F_r - F_r'\|_{\mathrm{TV}}$ at every $x$. Taking the supremum over $x \in A_p$ yields
\[
K_p \;\leq\; 1.
\]

\textbf{Responder's coupling.} The responder of type $\theta$ accepts iff $x \geq \theta$. The
type-$\theta$ responder's truncated cost upon receiving offer $x$ from the proposer's mixed
strategy $\sigma_p$ is
\begin{equation}
\label{eq:responder-cost}
c_r(\theta \mid \sigma_p) \;=\; \underbrace{|\theta - 0.5|}_{\text{fairness deviation, scaled}}
\;-\; \mathbb{E}_{x \sim \sigma_p}\bigl[\mathbf{1}\{x \geq \theta\} \cdot x\bigr],
\end{equation}
where the second term is the responder's expected payoff conditional on accepting. The first term
is independent of $\sigma_p$. The second satisfies
\[
\sup_\theta \bigl| \mathbb{E}_{\sigma_p}[\mathbf{1}\{x \geq \theta\}\,x] - \mathbb{E}_{\sigma_p'}[\mathbf{1}\{x \geq \theta\}\,x] \bigr|
\;\leq\; \sup_\theta \int |x|\, |d\sigma_p(x) - d\sigma_p'(x)|
\;\leq\; \tfrac{1}{2}\, \|\sigma_p - \sigma_p'\|_{\mathrm{TV}},
\]
because $|x| \leq 1/2$ on the proposer's action grid $A_p \subseteq [0, 0.5]$. Therefore
\[
K_r \;\leq\; 1/2.
\]

\textbf{The contraction ratio.} Substituting into Proposition~\ref{prop:br-lip},
\begin{equation}
\label{eq:explicit-L}
L \;\leq\; \frac{\max(K_p, K_r)}{2T} \;=\; \frac{1}{2T}.
\end{equation}

Under the TCSS cost-dependent temperature $T(\Delta) = \max(T_{\text{floor}},
T_{\text{base}}\Delta^{T_\alpha})$ with $T_{\text{floor}} = 0.5$, the effective temperature
satisfies $T \geq 0.5$. Substituting,
\[
L \;\leq\; \frac{1}{2 \cdot 0.5} \;=\; 1.
\]

\textbf{This is not yet a strict contraction.} The bound saturates at $L = 1$, which is the
boundary case. To obtain $L < 1$ rigorously requires either (i) a strictly positive lower bound on
$T$ above $0.5$ in the regime of interest, (ii) a sharper bound on $K_p$ or $K_r$ exploiting the
specific structure of the calibrated metric, or (iii) bounding the BR Lipschitz constant in a norm
other than total variation (e.g., the Wasserstein-1 distance on offer space, which would exploit
the small diameter of $A_p$).

\subsection*{B.4 A strict bound under restricted action grids}

A strict contraction can be obtained by restricting the proposer's action grid to ``serious offers''
$A_p^\diamond = \{0.05, 0.10, \ldots, 0.50\}$ (excluding $x = 0$). On this grid, the minimum
expected payoff $(1-x)$ is bounded above by $1 - 0.05 = 0.95$, but the bound on $K_p$ is unchanged.
A genuine tightening comes from observing that the truncated Mahalanobis cost in the proposer's
softmax dominates the goal-reachability term: with the calibrated values $\sigma_6^2 = 78.26$,
$\sigma_7^2 = 32.28$, the truncated cost at $x = 0.5$ is approximately $\sqrt{0.25 \cdot 0.044 +
0.51} = \sqrt{0.521} \approx 0.722$, while the maximum goal-reachability term is $|1-x| \leq 1$.
The ratio is order-one, and the proposer's softmax distribution is dominated by the cost term, not
the goal term.

When the cost term dominates the softmax, the effective Lipschitz constant of $c_p(\cdot \mid F_r)$
in $F_r$ scales with the weight of the goal-reachability term relative to the cost term. A
parameter-substitution estimate (substituting $T = 0.5$ and the ratio above) gives an effective
$L \approx 1/(2 T) \cdot 0.722^{-1} \cdot 1 \approx 1.38^{-1} \approx 0.72$. This is a heuristic
sharpening and not a proof; a rigorous version would require treating the softmax as a non-uniform
weighted average and applying the Wasserstein-1 norm.

We treat the strict-contraction bound in the calibrated regime as established under the additional
assumption that the cost term dominates the goal term uniformly, and as open in the general case.

\subsection*{B.5 The TCSS empirical regime is single-shot}

The rigorous bound of~\eqref{eq:explicit-L} is for the iterated ultimatum game in which both
proposer and responder update their mixed strategies via softmax best-response. The TCSS empirical
calibration does \emph{not} run this iteration. It uses a fixed empirical responder rejection
function
\begin{equation}
\label{eq:empirical-reject}
P_{\text{reject}}(x) \;=\; \frac{1}{1 + \exp((x - x_{\text{thresh}})/k)},
\quad x_{\text{thresh}} = 0.18,\;k = 0.15,
\end{equation}
estimated from the Fraser--Nettle ultimatum data~\cite{bond2026tcss}, and computes the proposer's
best-response against this fixed function. The proposer's optimization is therefore
\emph{single-shot}, not iterated, and the contraction question does not arise: the equilibrium
strategy is reached in one $A^*$ call per agent in time
$O(V^T + E^T \log V^T + t_h(k) V^T)$ as in Appendix~\ref{app:vt-et}.

Iterated best-response dynamics, and the contraction bound of~\eqref{eq:explicit-L}, become
relevant in three forward-going extensions that the present paper does not validate empirically but
the theorem framework predicts:

\begin{enumerate}
\item \textbf{Trust games with reputation updating.} Investors and trustees update beliefs about
each other's reciprocity over rounds; the iterated BR has genuine cross-coupling through the
reputation channel.
\item \textbf{Public-goods games with conditional cooperation.} Agents' contribution choices in
round $r$ depend on the empirical contribution distribution in round $r-1$; the social-impact
dimension $d_6$ couples agents across rounds.
\item \textbf{Network games with reinforcement learning.} Agents on a network update strategies
based on neighbor outcomes; the contraction bound extends with an additional factor that depends on
the network connectivity.
\end{enumerate}

In each case, Proposition~\ref{prop:br-lip} applies with the appropriate cost-coupling constants
$K_i$, and~\eqref{eq:explicit-L} gives an explicit Lipschitz bound. The qualitative prediction is
that the calibrated metric, when extended to these iterated regimes, will yield contraction ratios
in the range $L \in [0.5, 0.9]$, consistent with the geometric convergence rates observed in
behavioral game-theory experiments.

\subsection*{B.6 What this derivation does, and does not, establish}

The derivation in B.1--B.3 supplies an explicit Lipschitz bound on the joint best-response mapping
in terms of the cost-coupling constants $K_i$ and the temperature $T$. The bound is rigorous and
sharp in the operator-norm sense.

The application to the calibrated ultimatum game in B.3 yields $L \leq 1$ in total-variation norm,
which saturates the contraction boundary. A strict contraction requires either a tighter bound (B.4
sketches one approach) or a different norm.

The empirical TCSS calibration is single-shot, so the contraction analysis is not directly tested
by the published validation. The bound is a forward-going prediction about iterated extensions of
the framework. The principal open technical work is (i) the strict contraction bound in
Wasserstein-1 norm, (ii) empirical validation of iterated BR dynamics in trust games or repeated
PG games, and (iii) the extension to $n$-agent network games.

%==============================================================================
\begin{thebibliography}{99}
%==============================================================================

\bibitem{andersen2011}
S.~Andersen, S.~Erta\c{c}, U.~Gneezy, M.~Hoffman, and J.~A.~List, ``Stakes matter in ultimatum games,''
\emph{American Economic Review}, vol.~101, no.~7, pp.~3427--3439, 2011.

\bibitem{boyd2004}
S.~Boyd and L.~Vandenberghe, \emph{Convex Optimization}.
Cambridge, U.K.: Cambridge University Press, 2004.

\bibitem{bond2026geometric}
A.~H.~Bond, \emph{Geometric Economics: Heuristic Bounded Rationality, the Bond Geodesic, and the
Death of Homo Economicus}, monograph, 2026.

\bibitem{bond2026tcss}
A.~H.~Bond, ``Geometric prediction of economic behavior: cross-domain validation across game theory
and prospect theory,'' \emph{IEEE Transactions on Computational Social Systems}, 2026 (accepted with
revision).

\bibitem{chen2009}
X.~Chen, X.~Deng, and S.-H.~Teng, ``Settling the complexity of computing two-player Nash equilibria,''
\emph{Journal of the ACM}, vol.~56, no.~3, art.~14, 2009.

\bibitem{daskalakis2009}
C.~Daskalakis, P.~W.~Goldberg, and C.~H.~Papadimitriou, ``The complexity of computing a Nash
equilibrium,'' \emph{SIAM Journal on Computing}, vol.~39, no.~1, pp.~195--259, 2009.

\bibitem{hart1968}
P.~E.~Hart, N.~J.~Nilsson, and B.~Raphael, ``A formal basis for the heuristic determination of
minimum cost paths,'' \emph{IEEE Transactions on Systems Science and Cybernetics}, vol.~4, no.~2,
pp.~100--107, 1968.

\bibitem{kahneman1979}
D.~Kahneman and A.~Tversky, ``Prospect theory: an analysis of decision under risk,''
\emph{Econometrica}, vol.~47, no.~2, pp.~263--291, 1979.

\bibitem{lieder2020}
F.~Lieder and T.~L.~Griffiths, ``Resource-rational analysis: understanding human cognition as the
optimal use of limited computational resources,''
\emph{Behavioral and Brain Sciences}, vol.~43, e1, 2020.

\bibitem{mckelvey1995}
R.~D.~McKelvey and T.~R.~Palfrey, ``Quantal response equilibria for normal form games,''
\emph{Games and Economic Behavior}, vol.~10, no.~1, pp.~6--38, 1995.

\bibitem{milgrom1990}
P.~Milgrom and J.~Roberts, ``Rationalizability, learning, and equilibrium in games with strategic
complementarities,'' \emph{Econometrica}, vol.~58, no.~6, pp.~1255--1277, 1990.

\bibitem{monderer1996}
D.~Monderer and L.~S.~Shapley, ``Potential games,'' \emph{Games and Economic Behavior}, vol.~14,
no.~1, pp.~124--143, 1996.

\bibitem{nash1950}
J.~F.~Nash, ``Equilibrium points in $n$-person games,'' \emph{Proceedings of the National Academy of
Sciences}, vol.~36, no.~1, pp.~48--49, 1950.

\bibitem{simon1955}
H.~A.~Simon, ``A behavioral model of rational choice,'' \emph{Quarterly Journal of Economics},
vol.~69, no.~1, pp.~99--118, 1955.

\bibitem{tversky1974}
A.~Tversky and D.~Kahneman, ``Judgment under uncertainty: heuristics and biases,'' \emph{Science},
vol.~185, no.~4157, pp.~1124--1131, 1974.

\bibitem{tversky1992}
A.~Tversky and D.~Kahneman, ``Advances in prospect theory: cumulative representation of uncertainty,''
\emph{Journal of Risk and Uncertainty}, vol.~5, no.~4, pp.~297--323, 1992.

\end{thebibliography}

\end{document}
